\documentclass[UTF8,a4paper,12pt]{ctexart}

\usepackage{amsmath,amssymb,bm}
\usepackage{booktabs}
\usepackage{geometry}
\usepackage{longtable}
\usepackage{array}
\usepackage{enumitem}
\usepackage{microtype}
\usepackage{hyperref}

\geometry{left=2.6cm,right=2.6cm,top=2.5cm,bottom=2.5cm}
\setlength{\parindent}{2em}
\setlength{\parskip}{0.35em}
\renewcommand{\arraystretch}{1.35}
\setlist{nosep}
\hypersetup{hidelinks}

\title{方法与案例研究部分}
\author{}
\date{}

\begin{document}

\maketitle

本文方法与案例研究部分按照“问题刻画--方法执行--机制验证”的顺序展开。第 3 章首先把龙阳路站这类大型多线换乘站的应急疏散过程表达为多来源客流在共享设施服务链上的动态加载问题，回答本文到底要刻画什么；第 4 章说明该问题如何由当前状态 A* 对照方法和 AA 到达时刻队列感知搜索在同一执行层下求解，回答代码怎样把模型变成可运行过程；第 5 章再用低负荷和高负荷场景、来源组分配、关键设施负荷、整体疏散指标、Pathfinder 趋势对照和敏感性分析验证，回答该机制是否发生、是否改善全站疏散以及在哪些条件下成立。三章之间的关系不是并列介绍算法和实验，而是从换乘站疏散问题出发，经过路径组织方法，最终回到疏散效果解释。

\section{3 多来源客流作用下的换乘站疏散服务链模型}

龙阳路站应急疏散的关键困难并不在于乘客是否存在通向出口的可达路径，而在于多线路站台客流、列车到达释放客流、站厅滞留客流和换乘客流会在同一组楼扶梯、换乘通道、闸机和出口前区上叠加。尤其是换乘客流，其服务链通常跨越多个线路空间，既可能与本线出站客流竞争垂向设施，也可能在目标线路站厅和出入口处与其他来源客流再次汇合。若路径组织只依据当前可见阻抗，部分客流在到达关键设施时可能面对已经由前序批次形成的排队，从而使全站尾部清空受局部共享设施控制。

基于上述认识，本章将大型多线换乘站的疏散过程抽象为多来源客流在共享设施服务链上的动态分配问题。这里的“服务链”强调乘客从站台、列车、站厅或换乘通道入口出发后，需要依次经过若干空间单元和具有服务能力的设施，才能最终离站。模型章的任务不是介绍 AA，也不是单独证明某个搜索算法优越，而是定义研究对象、客流来源、设施容量、队列形成和评价指标。第 4 章中的路径搜索和共享容量执行均服务于本章提出的问题表达。

\subsection{3.1 多线换乘站服务链网络与客流来源}

将车站公共区抽象为有向网络 \(G=(V,E)\)。节点集合 \(V\) 表示站台候车区、列车释放点、站厅区域、换乘通道连接点、闸机区、出口前区和安全出口等空间单元；边集合 \(E\) 表示乘客可行走、可上行或可通过的连接关系。与一般最短路网络不同，本文并不把每条边都看成独立通行设施。对于楼梯、扶梯、换乘通道、闸机和出口前区等具有明确服务能力且可能被多条路径共同使用的对象，定义共享设施资源集合 \(R\)。若边 \(e\in E\) 对应某一物理共享设施，则通过映射
\[
\rho:E\rightarrow R\cup\varnothing
\]
把该边归入相应资源；若该边仅表示普通空间连接，则 \(\rho(e)=\varnothing\)。这样，同一楼梯或闸机即使在网络中被多条路径访问，也只对应一个服务容量，避免在路径计算中重复获得通行能力。

客流来源集合记为
\[
\mathcal G=\mathcal G^{p}\cup \mathcal G^{train}\cup \mathcal G^{hall}\cup \mathcal G^{tr},
\]
其中，\(\mathcal G^{p}\) 为站台等待或普通出站客流，\(\mathcal G^{train}\) 为列车到达释放客流，\(\mathcal G^{hall}\) 为站厅滞留客流，\(\mathcal G^{tr}\) 为换乘客流。本文把换乘客流作为解释整体疏散机制的关键来源组，而不是把研究目标缩小为换乘客流单独疏散。换乘客流在站内通常具有更长的服务链和更多跨层、跨线设施暴露，因此它会改变共享设施的到达负荷结构，并影响其他来源客流的排队和尾部离站。

离散仿真时间集合为
\[
\mathcal T=\{0,\Delta t,2\Delta t,\ldots,T\}.
\]
在任一决策时刻 \(t\)，可参与路径组织的客流批次集合记为 \(\mathcal B(t)\)。批次 \(b\in\mathcal B(t)\) 表示为
\[
b=(g_b,o_b,D_b,n_b,t_b),
\]
其中，\(g_b\in\mathcal G\) 为来源组，\(o_b\in V\) 为批次当前所在节点或队列入口，\(D_b\subseteq V\) 为可接受安全出口集合，\(n_b\) 为批次人数，\(t_b\) 为批次进入路径组织阶段的时刻。批次建模用于保留列车释放、站台等待、站厅滞留和换乘到达之间的时间差异，也便于在路径选择后记录该批次对共享设施的未来到达需求。

\begin{table}[htbp]
\centering
\caption{多来源客流与模型表达}
\small
\begin{tabular}{p{0.20\textwidth}p{0.32\textwidth}p{0.36\textwidth}}
\toprule
来源类型 & 模型表达 & 对整体疏散的作用 \\
\midrule
站台等待客流 & 由各线路站台等待区或分区节点释放 & 构成本线出站基础需求，并与列车释放流共同使用垂向设施 \\
列车到达释放客流 & 按线路、列车和车厢位置生成批次 & 放大短时间内的到达峰值，是高负荷场景中共享设施竞争的主要触发因素 \\
站厅滞留客流 & 由站厅 staging 节点进入网络 & 影响闸机、出口前区和部分站厅连接节点的空间接收状态 \\
换乘客流 & 由换乘通道、换乘楼扶梯或跨线入口节点生成 & 形成跨线服务链，与多线路出站流共同竞争楼扶梯、通道、闸机和出口前区 \\
\bottomrule
\end{tabular}
\end{table}

对批次 \(b\)，设 \(\mathcal P_b(t)\) 为从 \(o_b\) 到 \(D_b\) 中任一出口的候选可行路径集合。路径选择变量定义为
\[
x_{bp}(t)=
\begin{cases}
1, & \text{时刻 }t\text{ 批次 }b\text{ 选择路径 }p,\\
0, & \text{否则}.
\end{cases}
\]
在路径组织层，\(x_{bp}(t)\) 表示路径意图；在执行层，乘客是否能够按该意图前进，还要受共享设施容量、下游空间接收能力和在途移动过程约束。这一区分是本文模型的基础，因为路径规划结果不能直接等同于真实疏散结果。

\subsection{3.2 共享设施容量约束与队列形成机制}

设共享设施 \(r\in R\) 的服务率为 \(\mu_r\)，时间步内可释放能力为
\[
C_r(t)=\left\lfloor \mu_r\Delta t+\xi_r(t)\right\rfloor,
\]
其中，\(\xi_r(t)\) 为离散容量余量项。设施入口等待队列记为 \(Q_r(t)\)，节点 \(v\in V\) 的瞬时占用人数记为 \(N_v(t)\)，最大接收能力记为 \(\bar N_v\)。如果多个来源组在同一时间步到达同一共享设施，它们必须共同竞争 \(C_r(t)\)，而不能因为路径来源不同而分别获得容量。

设 \(I_{brp}(t,\theta)\) 表示批次 \(b\) 在时刻 \(t\) 选择路径 \(p\) 后是否预计于区间 \([\theta,\theta+\Delta t)\) 到达资源 \(r\)。资源 \(r\) 在 \(\theta\) 时刻的预测到达需求可写为
\[
A_r(\theta)=
\sum_{t\leq \theta}\sum_{b\in\mathcal B(t)}
\sum_{p\in\mathcal P_b(t)}
n_bx_{bp}(t)I_{brp}(t,\theta).
\]
在执行层，资源队列随到达需求和服务释放更新：
\[
Q_r(t+\Delta t)=\max\{Q_r(t)+A_r(t)-C_r(t),0\}.
\]
这条约束的意义不是形式上让路径层和执行层“保持一致”，而是保证同一物理设施在同一时间步只释放一次真实容量。对龙阳路站这样的多线换乘站而言，换乘流、站台流、列车释放流和站厅流可能经由不同路径进入同一楼扶梯、闸机或出口前区；若不在统一容量下执行，就无法判断路径组织是否真的缓解了共享设施竞争。

除服务能力外，下游空间接收能力也会影响乘客能否继续前进。设 \(B_v(t)\) 为当前时间步预计进入节点 \(v\) 的人数，则需要满足
\[
N_v(t)+B_v(t)\leq \bar N_v,\quad v\in V.
\]
当该约束不能满足时，乘客不应被无条件推入下游空间，而应在上游节点、边末端或设施入口形成等待。该机制用于表达高负荷疏散中排队向上游扩散的过程，也是后续分析“局部瓶颈如何影响尾部清空”的基础。

系统人数守恒为
\[
\sum_{v\in V}N_v(t)+\sum_{e\in E}M_e(t)+\sum_{r\in R}Q_r(t)+E^{out}(t)=N^{tot},
\]
其中，\(M_e(t)\) 为边上在途人数，\(E^{out}(t)\) 为累计离站人数，\(N^{tot}\) 为场景总人数。该式保证节点停留、边上移动、设施排队和完成疏散共同构成全系统状态，避免出现人数丢失或重复计数。

\subsection{3.3 到达时刻队列感知的路径组织目标}

多线换乘站疏散路径组织的难点并不是“有无路径”，而是同一条路径在不同到达时刻面对的设施服务状态可能完全不同。几何距离较短的路径，如果在乘客到达时已经由前序批次形成长队，则可能导致更长的实际疏散时间；几何距离较长的路径，如果避开未来队列峰值，则可能降低尾部清空时间和累计停滞等待。因此，本文把路径组织目标定义为到达时刻队列感知的广义代价最小化问题。

设当前可观测状态为
\[
\Omega(t)=\{N_v(t),Q_r(t),M_e(t),E^{out}(t),\mathcal A_r(t)\},
\]
其中，\(\mathcal A_r(t)\) 表示已登记的未来到达事件，包括在途客流和当前决策轮次内前序批次对资源 \(r\) 的预计到达。对候选到达时刻 \(\theta\geq t\)，资源 \(r\) 的预测队列可写为
\[
\widehat Q_r(\theta|t)=
\max\left\{
Q_r(t)+
\sum_{(\eta_j,a_j)\in\mathcal A_r(t),\eta_j\leq \theta}a_j
-\mu_r(\theta-t),0
\right\}.
\]
对应的预测等待时间为
\[
\widehat W_r(\theta|t)=\frac{\widehat Q_r(\theta|t)}{\mu_r}.
\]
对于批次自身进入共享设施后造成的平均服务占用，采用
\[
S_r(n_b)=\frac{n_b}{2\mu_r}.
\]
该项把批次规模转化为设施服务过程中的时间代价，而不是额外的人为惩罚。

给定候选路径 \(p=(e_1,e_2,\ldots,e_m)\)，设批次沿路径递推到达边 \(e_i\) 或其对应资源的时刻为 \(\eta_i\)，边 \(e_i\) 的行走时间为 \(\tau_{e_i}(\eta_i)\)，则路径预测广义代价表示为
\[
\widehat C_{bp}(t)=
\sum_{i=1}^{m}
\left[
\tau_{e_i}(\eta_i)+
I_{\rho(e_i)\in R}
\left(
\widehat W_{\rho(e_i)}(\eta_i|t)+S_{\rho(e_i)}(n_b)
\right)
\right]
+\lambda R_{bp}(t).
\]
其中，\(R_{bp}(t)\) 表示拥挤暴露或风险辅助项，\(\lambda\) 为尺度权重。路径组织的滚动子问题可写为
\[
p_b^{*}(t)=\arg\min_{p\in\mathcal P_b(t)}\widehat C_{bp}(t).
\]
这一目标并不追求简单最短距离，而是在行走时间、到达时刻等待、批次服务占用和拥挤暴露之间进行折中。第 4 章中的 AA 时间依赖搜索正是这一目标的求解方式。

\subsection{3.4 疏散效果评价指标}

为避免把研究写成单纯的算法快慢比较，本文从整体疏散效率、尾部清空、设施负荷分布、来源组差异和方法代价五个层面评价疏散效果。整体效率指标包括 \(T_{50}\)、\(T_{80}\)、\(T_{95}\)、\(T_{99}\)、\(T_{100}\)、平均疏散时间和平均站点通过率，其中
\[
T_{\theta}=\min\{t\mid E^{out}(t)\geq \theta N^{tot}\},\quad
\theta\in\{0.50,0.80,0.95,0.99,1.00\}.
\]
\(T_{95}\)、\(T_{99}\) 和 \(T_{100}\) 用于刻画大部分乘客离站后的尾部阶段，避免平均值掩盖少数滞留客流对安全组织的影响。

停滞与排队暴露指标包括累计静止人秒、资源队列人秒、空间阻塞人秒和关键设施峰值队列。累计资源队列人秒记为
\[
\Phi_Q=\sum_{t\in\mathcal T}\sum_{r\in R}Q_r(t)\Delta t,
\]
空间阻塞暴露记为
\[
\Phi_S=\sum_{t\in\mathcal T}\sum_{v\in V}S_v(t)\Delta t.
\]
若某一设施在峰值队列、累计队列人秒和持续时间上均较高，则说明其不是短时波动，而是对疏散过程产生连续约束的关键设施。

来源组指标用于解释换乘客流在整体疏散中的机制作用。对来源组 \(g\)，统计其离站曲线 \(E_g(t)\)、清空时间 \(T_g\)、关键设施暴露和最终出口分布。若换乘客流与列车释放客流在同一共享设施上形成到达叠加，则应进一步分析该设施是否控制全站尾部清空。

均衡性指标用于刻画出口和关键设施负荷分布。设 \(y_i\) 为出口或关键设施 \(i\) 的通过量，数量为 \(m\)，则 Jain 指数为
\[
J=\frac{(\sum_{i=1}^{m}y_i)^2}{m\sum_{i=1}^{m}y_i^2}.
\]
\(J\) 越接近 1，表示负荷分布越均衡。均衡性指标不能单独作为最优标准，但可以帮助判断路径组织是否把客流从少数拥挤设施转移到更多可用设施。

方法代价指标包括总移动距离、平均移动时间、有效疏散速度和墙钟计算时间。若 AA 在高负荷场景中降低尾部清空和累计停滞，但增加移动距离或计算时间，应如实写为效率改善与绕行、计算代价之间的权衡，而不是写成无条件优越。

\section{4 队列感知路径组织与共享容量执行方法}

第 3 章定义了多来源客流在共享设施服务链上的动态分配问题。本章进一步说明该问题如何在离散疏散仿真中被求解和执行。本文并不把路径搜索结果直接等同于疏散结果，而是将其视为当前时间步内的路径组织意图；所有意图仍需经过统一的资源服务能力、下游接收空间和在途到达过程执行。基于这一原则，本章首先说明客流批次和来源组的滚动生成方式，然后给出当前状态 A* 对照方法，再介绍 AA 的到达时刻队列预测与时间依赖搜索，最后说明共享容量执行层如何保证两种路径组织方法在同一物理约束下比较。

\subsection{4.1 客流批次生成与滚动决策流程}

在每个离散时间步开始时，系统读取上一时间步结束后的状态并冻结为当前决策状态 \(\Omega(t)\)。状态冻结的作用是使同一时间步内的路径组织从同一执行状态出发，避免由于批次处理顺序不同而随意改变节点占用、资源队列或下游接收能力。随后，系统识别当前可调度批次集合 \(\mathcal B(t)\)，包括已经位于节点等待区的乘客、到达设施入口但尚未获得服务的乘客，以及在当前时刻释放到站内网络中的列车到达客流。

当前程序中的客流输入由线路、来源类型和释放位置共同定义。基础客流包含各线路站台等待、站厅滞留和换乘来源组；高负荷场景在此基础上叠加多线路列车到达释放客流。初始化过程将客流分配到车厢、站台等待区、站厅 staging 节点和换乘入口节点，并保留来源组标识。这样，后续统计可以区分站台、列车、站厅和换乘来源组在出口利用、设施暴露和清空时间上的差异。

对每个批次，路径组织层根据当前状态和候选出口集合生成路径意图。若某一批次选择路径 \(p_b^{*}(t)\)，系统会记录其预计到达共享设施的时间、人数和来源组信息：
\[
\mathcal A_r(t)\leftarrow
\mathcal A_r(t)\cup\{(\widehat\eta_{brp},n_b,g_b,b)\}.
\]
该记录是本轮路径组织的未来到达承诺，使后续批次能够感知前序批次对同一设施容量的占用。需要强调的是，这只是路径组织层的预测信息，并不表示乘客已经通过该设施；实际通行仍由共享容量执行层决定。

一个时间步内的滚动决策流程包括六步：
\begin{enumerate}
\item 冻结当前节点占用、资源队列、在途客流和已登记未来到达事件；
\item 识别当前可参与路径组织的客流批次及其来源组；
\item 为批次生成候选出口和候选路径集合；
\item 调用当前状态 A* 对照方法或 AA 时间依赖搜索生成路径意图；
\item 把被接受路径上的预计资源到达写入未来到达事件集合；
\item 将全部路径意图交由共享容量执行层，更新下一时间步状态。
\end{enumerate}
该流程对应本文“路径组织--共享容量执行--状态更新”的基本求解框架。

\subsection{4.2 当前状态 A* 对照方法}

A* 是在加权图上求解最小代价路径的启发式搜索方法。对处于节点 \(v\) 的搜索状态，A* 使用评价函数
\[
f(v)=g(v)+h(v)
\]
确定优先扩展顺序。其中，\(g(v)\) 为从起点到节点 \(v\) 的累计代价，\(h(v)\) 为从节点 \(v\) 到目标节点的剩余代价估计。若 \(h(v)\) 是真实剩余代价的下界，A* 能够在减少无效扩展的同时保持路径搜索的合理性。

在静态路径规划中，边代价通常由距离、自由流时间或固定通行阻抗给出，同一个节点只需保留当前最小的 \(g(v)\)。但在地铁站疏散中，乘客到达楼扶梯、闸机或出口前区的时刻会影响其面对的队列长度和等待时间。同一节点在不同时刻被到达时，后续边代价可能不同。因此，基础 A* 可以作为路径搜索骨架，但不能完整表达到达时刻设施竞争。

本文使用 ImprovedAStar 作为计算对照方法。该方法在相同站内网络、相同客流输入、相同速度模型、相同共享设施容量、相同下游接收约束和相同终止条件下运行，用于提供当前状态路径组织的比较基准。它的作用是说明：如果路径评价主要依据当前可见状态和即时边权，而不显式登记本轮前序批次的未来到达负荷，疏散执行层会形成怎样的出口利用、关键设施排队和尾部清空结果。ImprovedAStar 仅承担算法对照角色，不对应已核验的现场疏散方案。

为了与 AA 保持可比，当前状态 A* 对照方法同样不直接决定乘客最终疏散结果。它输出的路径也要进入共享容量执行层，接受同一设施服务能力和下游接收能力约束。这样，实验比较只把差异限制在路径组织层，而不是把网络、客流、容量或统计口径混在一起。

\subsection{4.3 AA 时间依赖队列感知搜索}

本文将到达时刻队列感知的路径搜索方法记为 AA，即 AdaptiveQueueAwareAStar。AA 并不是重新定义第 3 章的模型，而是在 A* 搜索框架中引入到达时刻队列预测，使路径代价能够反映批次真正到达共享设施时可能面对的等待。其核心扩展包括三点：状态由单一节点扩展为时间标签，边代价由固定阻抗扩展为时间依赖代价，本轮已承诺的未来到达事件会影响后续批次的路径评价。

AA 对批次 \(b\) 从当前节点 \(o_b\) 到目标出口集合 \(D_b\) 进行多标签搜索。搜索标签表示为
\[
\ell=(v,\eta,c,\psi,\pi),
\]
其中，\(v\) 为标签所在节点，\(\eta\) 为预计到达该节点的时刻，\(c\) 为从起点到当前节点的累计广义代价，\(\psi\) 为累计拥挤暴露或风险代价，\(\pi\) 为前驱标签。同一节点上的两个标签即使位置相同，只要到达时刻不同，也可能面对不同的未来队列，因此不能简单按节点合并。AA 保留必要的时间标签，并通过支配关系删除明显劣势标签：若两个标签位于同一节点，且一个标签到达不晚、累计代价和拥挤暴露均不高，则另一个标签可被剔除。

当标签沿边 \(e=(u,v)\) 扩展时，首先计算物理行走时间 \(\tau_e(\eta)\)。若边 \(e\) 不对应共享设施，则新标签主要由行走时间更新；若 \(\rho(e)=r\)，则在预计进入资源的时刻 \(\eta\) 查询预测队列 \(\widehat Q_r(\eta|t)\)，并加入预测等待 \(\widehat W_r(\eta|t)\) 和批次自身服务占用 \(S_r(n_b)\)。标签扩展可概括为
\[
c'=c+\tau_e(\eta)+
I_{\rho(e)\in R}
\left(
\widehat W_{\rho(e)}(\eta|t)+S_{\rho(e)}(n_b)
\right)
+\lambda R_e(\eta).
\]
搜索使用自由流条件下到可用出口集合的最短剩余时间作为启发式下界：
\[
h(v)=\frac{d^{free}(v,D_b)}{v^{free}}.
\]
该下界只用于限制搜索空间，不替代真实路径代价。真实路径代价仍由已递推的行走时间、预测设施等待、批次服务占用和拥挤暴露共同确定。

AA 输出的路径不仅包括节点序列，还包括该批次预计到达各共享设施的时间、人数和来源组信息。若该路径被接受，这些信息会登记到当前时间步的未来到达事件集合，使随后处理的批次能够感知同一设施即将承受的负荷。由此，AA 的“自适应”不是泛泛地根据拥堵调权，而是随前序批次的路径承诺更新到达时刻队列预测，再据此重新评价后续批次路径。

\subsection{4.4 共享容量执行与结果导出}

路径搜索得到的是路径组织意图，不能直接等同于乘客已经按该路径完成疏散。执行层按照资源服务率、下游接收能力和在途移动过程推进乘客，记录实际队列、实际通行、空间阻塞、出口通过量和来源组清空时间。若同一共享设施在当前时间步只能服务 \(C_r(t)\) 人，则所有到达该设施入口的来源组共同竞争这一容量；若下游节点已达到接收上限，乘客即使获得服务也不能无限进入下游空间，而应在上游形成等待。

共享容量执行层的意义在于把第 3 章的设施竞争问题落到真实约束下检验。若路径组织阶段对某设施的未来负荷判断准确，执行层中该设施的队列峰值和持续时间应相应变化；若预测与执行发生偏差，下一时间步会根据更新后的状态重新组织路径。因此，本文比较 AA 与 ImprovedAStar 时，不能只看搜索得到的路径长度，而必须看路径进入执行层后的实际疏散曲线、设施队列和出口使用结果。

执行层输出包括总体疏散指标、线路清空时间、来源组出口分布、关键设施通过量、候选路径代价、负荷均衡指标和 Pathfinder 路径分配文件。这些输出决定第 5 章的实验顺序：先分析换乘及跨线客流如何被分配到不同出口和设施，再分析这种分配是否降低停滞等待和尾部清空，最后用 Pathfinder 趋势对照和敏感性分析检查结论边界。

\section{5 龙阳路站案例研究与机制验证}

本章以龙阳路站为案例，检验第 3 章提出的多来源客流共享设施服务链问题，以及第 4 章 AA 到达时刻队列感知路径组织方法的执行效果。案例研究的目的不是单纯比较两个算法名称下的总疏散时间，而是检验队列感知路径组织是否能够改变大型换乘站高负荷疏散中的客流加载方式，并由此影响全站尾部清空和等待暴露。为避免把偶然的总时间差异误读为机制有效，本文先分析换乘及跨线客流的出口和设施分配，再分析总体疏散时间、线路清空和关键设施负荷，最后通过 Pathfinder 微观趋势对照和敏感性分析讨论结论边界。

\subsection{5.1 案例站、客流输入与场景设置}

龙阳路站为多线换乘车站，站内包含多个线路站台、换乘通道、站厅、楼扶梯、闸机区和地面出口。案例网络根据 CAD 图纸和当前节点边关系建立；设施宽度、长度、连接关系和出口位置应在参数表中区分 CAD 可量测值、模型设定值、文献取值和必要假设值。尚未完成来源核验的设施容量和行为参数，不应写成实测值或官方值。速度--密度关系仅作为基础运动时间估计的一部分，本文实验不把它作为主要敏感性分析对象。

案例设置低负荷基础疏散场景和高负荷列车叠加疏散场景。低负荷场景包含站台等待、站厅滞留和换乘客流，用于观察共享设施竞争较弱时队列感知机制是否稳定；高负荷场景在基础客流上叠加多线路列车释放客流，用于放大多来源客流的到达重叠和共享设施容量竞争。根据当前模型配置，低负荷场景总人数为 2187 人，高负荷场景总人数为 17905 人；若后续仿真输入冻结时发生修改，正式结果表应同步更新。

\begin{table}[htbp]
\centering
\caption{案例场景设置}
\small
\begin{tabular}{p{0.23\textwidth}p{0.31\textwidth}p{0.34\textwidth}}
\toprule
场景 & 客流组成 & 验证目的 \\
\midrule
低负荷基础疏散场景 & 站台等待客流、站厅滞留客流和换乘客流；当前模型配置总人数为 2187 人 & 检查共享设施竞争较弱时，AA 是否保持稳定路径组织，避免为预测队列引入无必要绕行或额外等待 \\
高负荷列车叠加疏散场景 & 基础客流叠加多线路列车到达释放客流；当前模型配置总人数为 17905 人 & 放大多来源客流到达重叠，检验到达时刻队列预测对关键设施负荷、等待暴露和尾部清空的作用 \\
\bottomrule
\end{tabular}
\end{table}

两种路径组织方法采用相同站内网络、相同客流输入、相同设施容量、相同下游接收约束和相同统计指标。ImprovedAStar 作为当前状态 A* 计算对照方法，AA 作为到达时刻队列感知方法。这样设置的目的是把差异集中在路径组织层，便于解释 AA 是否通过改变多来源客流的设施加载方式影响整体疏散。

\subsection{5.2 换乘与跨线客流分配结果分析}

本节先分析客流分配，而不是直接报告总疏散时间。这样安排是因为本文的核心结论应由“路径组织如何改变共享设施负荷”支撑。如果只比较 \(T_{100}\)，读者无法判断改善来自换乘流分散、出口利用均衡、关键楼扶梯排队降低，还是来自其他偶然因素。因此，实验首先比较两种方法下各线路、各来源组和跨线客流的最终出口分布，以及关键共享设施的通过量和队列暴露。

对每个来源组 \(g\)，统计其最终出口集合、跨线出口比例、经过关键共享设施的人数和离站曲线。跨线出口在这里不是评价某类乘客是否“被影响”的指标，而是判断路径组织是否改变了多线换乘站中服务链加载方式的证据。如果 AA 在高负荷场景中引导部分客流避开未来队列峰值较高的设施，并使用其他尚有余量的出口或通道，则应在出口分布和关键设施通过量上体现出来。

\begin{table}[htbp]
\centering
\caption{换乘与跨线客流分配结果记录表}
\small
\begin{tabular}{p{0.17\textwidth}p{0.16\textwidth}p{0.17\textwidth}p{0.20\textwidth}p{0.20\textwidth}}
\toprule
场景 & 方法 & 跨线客流规模 & 主要变化位置 & 解释重点 \\
\midrule
低负荷基础疏散 & ImprovedAStar & 待冻结 & 待冻结 & 当前状态路径组织下的基础出口分布 \\
低负荷基础疏散 & AA & 待冻结 & 待冻结 & 检查非强瓶颈条件下是否出现过度分流 \\
高负荷列车叠加 & ImprovedAStar & 待冻结 & 待冻结 & 识别多来源客流是否集中加载少数共享设施 \\
高负荷列车叠加 & AA & 待冻结 & 待冻结 & 判断到达时刻队列预测是否改变跨线与换乘相关服务链分配 \\
\bottomrule
\end{tabular}
\end{table}

结果解释应围绕三个问题展开。第一，换乘流是否与列车释放流共同进入同一楼扶梯、通道、闸机或出口前区；第二，AA 是否降低了少数设施的集中加载，并提高出口或关键设施负荷均衡性；第三，这种分配变化是否为后续尾部清空和停滞等待变化提供机制解释。只有当分配证据和效率证据方向一致时，才能把 AA 的改善解释为队列感知路径组织的作用。

\subsection{5.3 全站疏散效率与瓶颈转移分析}

在完成分配机制分析后，本节比较整体疏散效率和关键瓶颈变化。整体效率指标包括 \(T_{50}\)、\(T_{80}\)、\(T_{95}\)、\(T_{99}\)、\(T_{100}\)、平均疏散时间、累计静止人秒和平均站点通过率。低负荷场景的重点不是追求显著差异，而是检查两种方法在瓶颈较弱时是否得到接近的总清空结果；高负荷场景的重点是检验 AA 是否降低尾部清空、累计停滞和关键设施排队。

下表中的“待填”表示相应结果在最终仿真输出冻结后填入。

\begin{table}[htbp]
\centering
\caption{整体疏散效率比较}
\small
\begin{tabular}{p{0.15\textwidth}p{0.15\textwidth}p{0.12\textwidth}p{0.12\textwidth}p{0.12\textwidth}p{0.12\textwidth}p{0.16\textwidth}}
\toprule
场景 & 方法 & \(T_{50}\) & \(T_{95}\) & \(T_{99}\) & \(T_{100}\) & 解释重点 \\
\midrule
低负荷 & ImprovedAStar & 待填 & 待填 & 待填 & 待填 & 计算对照基准 \\
低负荷 & AA & 待填 & 待填 & 待填 & 待填 & 检查稳定性和额外代价 \\
高负荷 & ImprovedAStar & 待填 & 待填 & 待填 & 待填 & 当前状态路径组织下的尾部表现 \\
高负荷 & AA & 待填 & 待填 & 待填 & 待填 & 到达时刻队列预测下的尾部变化 \\
\bottomrule
\end{tabular}
\end{table}

除总疏散时间外，还应比较队列暴露和方法代价。当前最新未冻结输出显示，高负荷场景中 AA 的尾部清空和累计静止人秒呈下降趋势，同时总移动距离和墙钟计算时间增加。正式稿在结果冻结后应把该现象写成“等待暴露降低与绕行、计算代价之间的权衡”，而不是写成 AA 在所有指标上均优。低负荷场景中，若 \(T_{100}\) 接近或相同，应解释为瓶颈不强时队列预测机制的边际收益有限，这恰好说明低负荷场景用于稳定性和边界检验，而不是用来证明显著提效。

\begin{table}[htbp]
\centering
\caption{瓶颈暴露与方法代价比较}
\small
\begin{tabular}{p{0.14\textwidth}p{0.15\textwidth}p{0.14\textwidth}p{0.15\textwidth}p{0.15\textwidth}p{0.13\textwidth}p{0.12\textwidth}}
\toprule
场景 & 方法 & 累计静止人秒 & 关键设施 Jain 指数 & 出口 Jain 指数 & 总移动距离 & 计算耗时 \\
\midrule
低负荷 & ImprovedAStar & 待冻结 & 待冻结 & 待冻结 & 待冻结 & 待冻结 \\
低负荷 & AA & 待冻结 & 待冻结 & 待冻结 & 待冻结 & 待冻结 \\
高负荷 & ImprovedAStar & 待冻结 & 待冻结 & 待冻结 & 待冻结 & 待冻结 \\
高负荷 & AA & 待冻结 & 待冻结 & 待冻结 & 待冻结 & 待冻结 \\
\bottomrule
\end{tabular}
\end{table}

线路清空结果用于进一步判断尾部滞留是否集中在特定线路或来源组。若高负荷场景中最后清空线路在两种方法下保持一致，但 AA 缩短了该线路尾部时间，则说明方法主要缓解了最后阶段的排队释放；若最后清空线路发生变化，则应分析瓶颈是否从原有设施转移到新的共享设施，避免把瓶颈转移误写成瓶颈消除。

\subsection{5.4 Pathfinder 微观对照与敏感性分析}

Pathfinder 对照用于检查图模型结论在微观行人仿真中的趋势表现。对照内容包括疏散曲线、拥堵位置、出口利用和来源组滞留特征。由于 Pathfinder 的个体避让、局部冲突和空间占用机制与图模型不同，本文不应把它写成逐个乘客轨迹的严格校准，而应说明两类模型是否在关键瓶颈位置、拥堵持续时段和总体疏散曲线形态上保持一致。

\begin{table}[htbp]
\centering
\caption{图模型与 Pathfinder 对照维度}
\small
\begin{tabular}{p{0.20\textwidth}p{0.34\textwidth}p{0.34\textwidth}}
\toprule
对照维度 & 图模型输出 & Pathfinder 输出 \\
\midrule
疏散曲线 & 累计疏散人数、\(T_{95}\)、\(T_{99}\)、\(T_{100}\) & 个体离站时间分布和累计疏散曲线 \\
拥堵位置 & 关键设施队列峰值、队列人秒和空间阻塞人秒 & 密度云图、局部聚集区域和拥堵持续时段 \\
出口利用 & 各出口通过人数及来源组构成 & 各出口实际离站人数和占比 \\
来源组差异 & 站台、列车、站厅和换乘客流清空时间 & 不同人群组离站时间和拥堵暴露 \\
\bottomrule
\end{tabular}
\end{table}

若 Pathfinder 中的高密度区域与图模型识别的高队列资源一致，且疏散曲线尾部变化方向相同，则可认为微观仿真在趋势上支持路径组织模型结论。若两者不一致，应优先检查空间几何表达、局部避让机制、设施容量设定和出口边界条件，而不能直接把差异解释为某一方法失效。

敏感性分析应服务于本文机制，而不是泛泛改变所有行人参数。本文建议重点分析三类因素：换乘客流比例、列车到达重叠程度和关键共享设施服务能力。换乘客流比例用于检验服务链交叉强度变化时 AA 的作用边界；列车到达重叠程度用于检验多来源客流在时间上集中释放时到达时刻预测的价值；关键共享设施服务能力用于检验瓶颈强弱变化对路径组织效果的影响。

\begin{table}[htbp]
\centering
\caption{敏感性分析设置与解释目的}
\small
\begin{tabular}{p{0.24\textwidth}p{0.30\textwidth}p{0.34\textwidth}}
\toprule
敏感性因素 & 调整方式 & 解释目的 \\
\midrule
换乘客流比例 & 在总需求或基础需求可控条件下调整换乘来源组占比 & 检验服务链交叉增强时，AA 是否更能识别共享设施未来排队 \\
列车到达重叠程度 & 调整多线路列车释放的时间间隔或同步程度 & 检验到达高峰越集中时，到达时刻队列预测是否更重要 \\
关键共享设施服务能力 & 调整代表性楼扶梯、通道、闸机或出口前区容量 & 检验瓶颈强弱变化下，AA 的适用范围和失效边界 \\
\bottomrule
\end{tabular}
\end{table}

结果解释应围绕适用边界展开。若 AA 的优势主要出现在高换乘比例、高到达重叠或关键设施服务能力较低的条件下，应将其解释为面向共享设施竞争明显场景的路径组织方法；若在低负荷或瓶颈较弱场景中差异较小，则说明模型不会在非瓶颈条件下强行产生显著改善。这样的解释比单纯报告“AA 更快”更能说明方法适用于哪些多线换乘站疏散情境。

\end{document}
